% Options for packages loaded elsewhere
\PassOptionsToPackage{unicode}{hyperref}
\PassOptionsToPackage{hyphens}{url}
\documentclass[
]{article}
\usepackage{xcolor}
\usepackage{amsmath,amssymb}
\setcounter{secnumdepth}{-\maxdimen} % remove section numbering
\usepackage{iftex}
\ifPDFTeX
  \usepackage[T1]{fontenc}
  \usepackage[utf8]{inputenc}
  \usepackage{textcomp} % provide euro and other symbols
\else % if luatex or xetex
  \usepackage{unicode-math} % this also loads fontspec
  \defaultfontfeatures{Scale=MatchLowercase}
  \defaultfontfeatures[\rmfamily]{Ligatures=TeX,Scale=1}
\fi
\usepackage{lmodern}
\ifPDFTeX\else
  % xetex/luatex font selection
\fi
% Use upquote if available, for straight quotes in verbatim environments
\IfFileExists{upquote.sty}{\usepackage{upquote}}{}
\IfFileExists{microtype.sty}{% use microtype if available
  \usepackage[]{microtype}
  \UseMicrotypeSet[protrusion]{basicmath} % disable protrusion for tt fonts
}{}
\makeatletter
\@ifundefined{KOMAClassName}{% if non-KOMA class
  \IfFileExists{parskip.sty}{%
    \usepackage{parskip}
  }{% else
    \setlength{\parindent}{0pt}
    \setlength{\parskip}{6pt plus 2pt minus 1pt}}
}{% if KOMA class
  \KOMAoptions{parskip=half}}
\makeatother
\setlength{\emergencystretch}{3em} % prevent overfull lines
\providecommand{\tightlist}{%
  \setlength{\itemsep}{0pt}\setlength{\parskip}{0pt}}
\usepackage{bookmark}
\IfFileExists{xurl.sty}{\usepackage{xurl}}{} % add URL line breaks if available
\urlstyle{same}
\hypersetup{
  hidelinks,
  pdfcreator={LaTeX via pandoc}}

\author{}
\date{}

\begin{document}

\textbf{RNR Coding --- Experimental Design Companion}

\emph{Empirical validation plan for the framework paper, v0.7}

Pavlo Ratushnyi --- May 2026

\textbf{1. Purpose and scope}

This document specifies the experimental program required to convert the
theoretical claims of RNR Coding v0.5 into empirically verified
statements. The main paper is theory-only by design: every quantitative
prediction in its §11 is labeled as a conjecture, every theorem is
conditional on stated assumptions, and Remark 5.7 invokes the manifold
hypothesis without verifying it on specific data. This companion
document fills the gap by stating exactly what must be measured, on what
data, with what baselines, under what statistical methodology, and how
the results should be reported to falsify or confirm each theoretical
claim.

The document covers three classes of empirical work: (a) coding-rate
measurements that test each Type's predicted advantage against
baselines; (b) information-theoretic measurements (total correlation
$\mathrm{TC}(D_N)$ and intrinsic dimension $d$) that test whether
specific natural data classes satisfy the conditions of Theorem 5.6
and Remark 5.7; and (c) cross-platform reproducibility tests that
verify Theorems 10.1 and 10.2 hold for a specific reference
implementation.

Scope deliberately excluded: privacy auditing (a separate effort with
its own threat model), end-user benchmarking (a deployment concern), and
large-scale model training (the framework assumes a small distilled
predictor; training methodology is upstream of this evaluation).

\textbf{1.1 Relation to the main paper}

This document does not modify any claim in the main paper. The main
paper\textquotesingle s §13.3 explicitly marks empirical validation as
out of scope; this document is the artifact that §13.3 points to.
Results from executing this plan would normally be published as a
separate experimental paper, jointly authored or independently. Both
documents stand alone.

\textbf{2. Hypotheses to test}

Each numbered hypothesis below corresponds to one or more theorems in
the main paper. Hypotheses are stated in falsifiable form with explicit
measurement procedures. The relationship to theory is given as either
confirming (the theorem holds; the experiment is checking that
conditions are realistic) or quantifying (the theorem\textquotesingle s
effect size on real data).

\textbf{2.1 Type-I rate (Theorems 4.1, 4.2)}

\textbf{H1 {[}quantifying{]}:} On English-language text and source code
corpora, an RNR Type-I coder with the Code12 $E_{12}$-A mapping and a small
(5--20 MB amortized) byte-level transformer predictor achieves coding
rate 1.5--2.0 bits per byte, competitive with LZMA at maximum preset
(2.1--2.5 bpb) and behind NNCP (1.3--1.6 bpb). The gap to NNCP narrows
under $E_{12}$-C (learned mapping).

\emph{Falsification:} if Type-I rate exceeds 2.5 bpb on enwik9 (worse
than LZMA), or if Type-I rate exceeds NNCP rate by more than 0.5 bpb
after $E_{12}$-C optimization, the conjecture in §11 is falsified for this
regime.

\textbf{2.2 Type-II separation on structured data (Theorems 5.3, 5.4,
5.6)}

\textbf{H2 {[}quantifying{]}:} On database pages, structured binary
records (Protocol Buffers, MessagePack, Apache Parquet), and instruction
streams (x86-64, RISC-V), an RNR Type-II coder with semi-non-factorized
positional model achieves coding rate 10--20\% below LZMA at maximum
preset.

\textbf{H3 {[}confirming{]}:} For each of the data classes in H2,
the measured total correlation $\mathrm{TC}(D_N)$ is $\Theta(N)$,
i.e., remains at least 10\% of $N$ as $N$ grows over the range 1
KiB to 16 KiB.

\emph{Falsification:} if measured $\mathrm{TC}(D_N)$ decays as
$o(N)$ for any of the data classes, then Theorem 5.6's separation
regime does not apply and Type-II provides only logarithmic-scale
gains. If Type-II rate equals or exceeds LZMA on database pages,
the separation regime is not realized at the predictor size tested.

\textbf{2.3 Manifold-hypothesis quantification (Remark 5.7)}

\textbf{H4 {[}confirming{]}:} For each natural-data class in H2 and
H6, the estimated intrinsic dimension $d$, computed via the
Levina-Bickel maximum-likelihood estimator on byte-level embedding
vectors at multiple ambient dimensions, lies in the regime
$d \ll N$.

\emph{Falsification:} if $d$ is comparable to $N$ (within a factor
of 2) for any tested natural-data class, the manifold-hypothesis
quantification in Remark 5.7 does not apply to that class, and
$\mathrm{TC}(D_N)$ is not predicted to be $\Theta(N)$ for it. The
conditional theorem still holds; only the empirical applicability is
in question.

\textbf{2.4 Type-III-A mode-selection effectiveness (Theorems 6.1,
6.2)}

\textbf{H5 {[}quantifying{]}:} On heterogeneous archives (mixed
text/binary, container files, database dumps with header and
payload), an RNR Type-III-A scheduler with calibrated $U_m(t)$
achieves rate within 5\% of the per-tile oracle best mode. The
per-mode prediction error $\varepsilon_t$ averages below 0.1 bits
per position.

\emph{Falsification:} if the scheduler's rate exceeds the per-tile
oracle by more than 20\% on heterogeneous data, the predictor for
$U$ is undercalibrated and the $2\varepsilon$ bound is not
approached in practice.

\textbf{2.5 Type-III-B mutual-information advantage (Theorem 6.3)}

\textbf{H6 {[}quantifying{]}:} Measured $I(H ; L \mid C)$ on ASCII
text falls in 0.3--0.7 bits per byte; on x86-64 instruction streams,
0.4--0.9 bits per byte; on float32 numeric arrays (scientific data),
1.0--2.0 bits per byte.

\emph{Falsification:} if $I(H ; L \mid C)$ is below 0.1 bits per
byte on any of these classes, the Type-III-B advantage over Type-II
at nibble alphabet is negligible and not worth the additional model
cost.

\textbf{2.6 Type-III-C with bits-back (Theorem 6.6)}

\textbf{H7 {[}confirming{]}:} An RNR Type-III-C coder with
bits-back ANS over a small VAE-style latent predictor achieves rate
within 2\% of $H(X \mid Y)$ on natural image data, matching the
bits-back identity. The empirical
$\mathrm{KL}(q \,\Vert\, P_{\text{posterior}})$ gap is bounded by
the variational training loss.

\emph{Falsification:} if the measured rate exceeds
$H(X \mid Y) + 5\%$ on image data, the bits-back accounting in
Theorem 6.6 has implementation defects or the variational posterior
$q$ is poorly fit. The theorem is unaffected.

\textbf{2.7 Reproducibility (Theorems 10.1, 10.2)}

\textbf{H8 {[}confirming{]}:} A reference RNR implementation built to
Definition 10.1 produces byte-identical archive outputs across (a)
x86-64 Linux with stock libc, (b) ARM64 Linux on Apple Silicon, (c)
RISC-V QEMU emulation, (d) Windows x86-64 with MSVC runtime, on a
100-block test corpus.

\emph{Falsification:} if any platform produces a byte-divergent archive
on any test block, the reference implementation does not meet Theorem
10.1\textquotesingle s conditions and must be fixed. The theorem is
unaffected.

\textbf{H9 {[}confirming{]}:} Online adaptation via a compatible
integer-arithmetic optimizer (integer SGD, signSGD, or integer Adam per
Definition 10.1) produces byte-identical archives across the same four
platforms as H8, after 1000 adaptation steps.

\textbf{2.8 Random access scaling (Theorems 10.3, 10.4)}

\textbf{H10 {[}confirming{]}:} For RNR archives built with a
$W$-bounded transformer predictor ($W = 4$ KiB) and sync points
placed every $K = 64$ KiB, wall-clock time for random access to
$k = 1$ KiB at any position $N$ satisfies: (a)~$T(N, k)$ is
independent of $N$ within measurement noise, across archive sizes
spanning 100 MiB to 100 GiB; (b)~$T(N, k)$ lies in the range 20 to
200 milliseconds on commodity hardware; (c)~ratio overhead from
sync points relative to a no-sync RNR baseline is below 2 percent
for $K = 64$ KiB and below 5 percent for $K = 4$ KiB.

\textbf{H11 {[}quantifying{]}:} For the LSTM-based RNR variant
(non-$W$-bounded, requiring state snapshots), random-access time
$T(N, k)$ is 3 to 10 times higher than the $W$-bounded transformer
variant at matched $K$, due to state-snapshot deserialization and
predictor warm-start cost. Seek index size grows by approximately
$L(\text{state})$ per sync point.

\emph{Falsification:} if $T(N, k)$ grows monotonically with $N$ for
the $W$-bounded variant, the bounded-context assumption is not
realized in the implementation, or sync points are ineffective. If
sync-point ratio overhead exceeds 10 percent at $K = 64$ KiB, the
per-sync-point cold-context cost is larger than predicted,
requiring reconsideration of $W$ or $K$.

\emph{Comparison baselines:} $T(N, k)$ for sequential decoding of
NNCP at the same archive sizes (expected: $O(N)$ scaling, hours for
1 GiB); $T(N, k)$ for bgzip and zstd seekable mode at matched
block sizes (expected: comparable to RNR $W$-bounded). The
kill-feature claim of §10.7 stands or falls on this comparison.

\textbf{2.9 Mounted-filesystem performance (Theorem 10.5)}

\textbf{H12 {[}confirming{]}:} A FUSE-based driver wrapping an RNR
archive of a typical Linux distribution ISO (5 to 20 GiB) supports
standard filesystem operations (open, read, stat, readdir) with
per-operation latency within 3 times the latency of the same operations
on a squashfs-compressed version of the same ISO. Sequential read
throughput of large files is within 2 times the squashfs throughput on
the same hardware. Compression ratio is at least 1.5 times better than
squashfs+zstd on the same ISO.

\textbf{H13 {[}confirming{]}:} Random 4 KiB reads issued via the mounted
filesystem at uniformly random offsets show 95th-percentile latency
below 100 milliseconds on CPU SIMD path, below 5 milliseconds on GPU
path, and below 1 millisecond on NPU path, for a small distilled
byte-level transformer predictor (10 to 50 million int8 parameters).

\emph{Falsification:} if mounted-filesystem latency exceeds 5 times
squashfs latency on any tested platform, the inference cost dominates
and the deployment is not interactive. If compression ratio is not
improved versus squashfs+zstd, RNR offers no deployment-time benefit
over the existing alternative.

\textbf{2.10 Hardware-target portability (Theorem 10.6)}

\textbf{H14 {[}confirming{]}:} A reference RNR predictor implementation,
ported to each of (a) x86-64 scalar CPU, (b) AVX-512 SIMD CPU, (c)
NVIDIA GPU using int8 Tensor Cores, (d) AMD GPU using matrix
accelerators, (e) Apple Neural Engine, (f) Google Edge TPU (if
accessible), produces byte-identical predictor output on a fixed
1000-input test set. Cross-target byte-identity rate is 100 percent.

\emph{Falsification:} any divergent output between two conforming
targets falsifies the hardware-target portability claim and indicates
that Theorem 10.6\textquotesingle s conformance specification is
incomplete for the diverging path. The fix is to identify the specific
integer-arithmetic or lookup-table behavior that diverges and either (i)
revise the conformance specification to require it explicitly, or (ii)
exclude the diverging target from the conforming-target set.

\textbf{3. Corpus selection}

The corpus is chosen to exercise each Type on its target regime, with
explicit negative controls. All corpora are public or publicly
reproducible to support independent verification.

\textbf{3.1 Text and source code}

\begin{itemize}
\item
  enwik9 (first $10^9$ bytes of English Wikipedia, standard Hutter
  Prize corpus). Tests Type-I and Type-III-B on ASCII text.
\item
  Linux kernel source tarball (latest LTS release, .tar without
  compression). Tests Type-I and Type-III-C on source code with
  structural repetition.
\item
  CommonCrawl HTML subset (10 GB, fixed snapshot hash). Tests Type-III-C
  with JSON/HTML token dictionaries.
\end{itemize}

\textbf{3.2 Structured binary}

\begin{itemize}
\item
  PostgreSQL pg\_dump output from a real schema (10 GB, anonymized or
  synthetic). Tests Type-II positional regularity.
\item
  Protocol Buffers serialized messages from a public schema (e.g.,
  TensorFlow saved-model files). Tests Type-III-C structural-marker
  tokens.
\item
  Apache Parquet column files. Tests Type-II on columnar layouts.
\end{itemize}

\textbf{3.3 Image and numeric data}

\begin{itemize}
\item
  CIFAR-10 raw uncompressed pixel data (50000 $32 \times 32$
  images, 150 MB). Tests Type-III-B and Type-III-C with bits-back on
  small images.
\item
  ImageNet validation set raw RGB (50000 images at original resolution).
  Tests Type-III-C with bits-back at scale.
\item
  Bayer-mosaic DSLR raw files (sample dataset of 1000 NEF or DNG files,
  public mirror). Tests Type-II on positional sensor structure.
\item
  HDF5 scientific arrays (sample subset of CERN or LSST public data).
  Tests Type-III-B on float32 numeric columns.
\end{itemize}

\textbf{3.4 Heterogeneous and negative controls}

\begin{itemize}
\item
  Linux distribution ISO (one specific release, archived for
  reproducibility). Tests Type-III-A heterogeneous scheduling.
\item
  64 MB of /dev/urandom output (negative control: incompressible).
\item
  32 MB of AES-CTR encrypted text (negative control: appears random).
\item
  128 MB of zstd-precompressed text (negative control: high-entropy
  residual).
\end{itemize}

All corpora are versioned by SHA-256 hash; experimental scripts
reference corpora by hash, not filename, to ensure long-term
reproducibility.

\textbf{4. Baseline implementations}

Every measurement compares the RNR variant against multiple baselines
drawn from classical, dictionary-based, and neural lossless compressors.
Baselines are pinned to specific releases.

\textbf{4.1 Classical and dictionary-based}

\begin{itemize}
\item
  gzip 1.10+ at level 9 (low baseline).
\item
  bzip2 1.0.8 at default settings (block-sorting baseline).
\item
  xz/LZMA 5.4+ at preset -9e (extreme preset, current ceiling for
  general-purpose dictionary compression).
\item
  zstd 1.5+ at level 22 (high-throughput modern baseline).
\item
  zpaq 7.15 at -m5 (long-range context mixing baseline).
\end{itemize}

\textbf{4.2 Context-mixing and neural}

\begin{itemize}
\item
  Cmix v19+ (current state-of-the-art among non-neural context-mixing
  coders).
\item
  NNCP latest release (Bellard, current state-of-the-art neural
  byte-level coder).
\item
  PixelCNN++ reference implementation (Salimans et al. 2017 release) on
  image corpora only.
\item
  Bit-Swap reference implementation (Kingma et al. 2019) on image
  corpora only.
\item
  IDF++ reference implementation (van den Berg et al. 2021) on image and
  numeric corpora.
\end{itemize}

\textbf{4.3 RNR implementations under test}

Six RNR variants are evaluated, plus their Type-III-A composition:

\begin{itemize}
\item
  RNR-I-A: Type-I with Code12 $E_{12}$-A (shortened Hamming) mapping.
\item
  RNR-I-C: Type-I with $E_{12}$-C (learned mapping, trained per corpus).
\item
  RNR-II: Type-II semi-non-factorized positional coder.
\item
  RNR-III-B: Type-III-B byte tensor coder.
\item
  RNR-III-C-D: Type-III-C with discrete BPE-style dictionary
  ($1024$ tokens, $32 \times 32$ grid).
\item
  RNR-III-C-V: Type-III-C with continuous latents via bits-back ANS over
  a VAE-style model.
\item
  RNR-Composed: Type-III-A scheduler over the above six.
\end{itemize}

All RNR variants share the same underlying byte-level transformer
predictor (8--16 MB amortized), with type-specific output heads. This
isolates the effect of the representation choice from the effect of
predictor capacity.

\textbf{5. Metrics}

\textbf{5.1 Primary: coding rate}

Bits per byte, measured as (codeword size in bits) / (uncompressed size
in bytes), reported separately for each (coder, corpus) pair. Report
includes:

\begin{itemize}
\item
  Total bits per byte (full archive).
\item
  Amortized rate excluding $L(M)$ (relevant for large archives where
  model cost vanishes).
\item
  Per-block rate distribution: minimum, 25th percentile, median, 75th
  percentile, maximum.
\end{itemize}

\textbf{5.2 Secondary: computational cost}

\begin{itemize}
\item
  Compression wall-clock time per MiB of input.
\item
  Compression peak resident memory (RSS).
\item
  Decompression wall-clock time per MiB of output.
\item
  Decompression peak resident memory (RSS).
\item
  CPU utilization (single-core; parallel speedup reported separately if
  multi-threaded).
\end{itemize}

\textbf{5.3 Information-theoretic measurements}

Independent of compression rate, several information-theoretic
quantities are estimated to test the conditions of Theorem 5.6 and
Remark 5.7:

\begin{itemize}
\item
  Total correlation $\mathrm{TC}(D_N)$ via the Cmix entropy
  estimator at multiple block lengths
  $N \in \{256, 512, 1024, 2048, 4096, 8192, 16384\}$. Estimated as
  (sum of marginal entropies measured from histograms) minus
  (Cmix-estimated joint entropy).
\item
  Intrinsic dimension $d$ via the Levina-Bickel maximum-likelihood
  estimator with $k \in \{5, 10, 20, 50\}$ nearest neighbors, computed
  on byte $n$-gram embedding spaces ($n \in \{2, 4, 8, 16\}$). Report
  $d$ as a function of $(k, n)$.
\item
  Conditional mutual information $I(H ; L \mid C)$ for Type-III-B,
  estimated via the empirical joint and marginal distributions of
  (high nibble, low nibble) conditioned on a sliding-window context.
\item
  Per-mode prediction error $\varepsilon_t$ for Type-III-A, computed
  as the empirical maximum over modes of $|U_m(t) - L_m(t)|$ over a
  sample of tiles.
\end{itemize}

\textbf{5.4 Reproducibility metrics}

\begin{itemize}
\item
  Byte-identity rate across (platform\_i, platform\_j) pairs for each
  test block, reported as (number of byte-identical archives) / (total
  archives).
\item
  Earliest divergent byte offset on first divergence, if any, for
  diagnostic purposes.
\item
  Floating-point fallback count: any tile that, under the Theorem 10.1
  specification, would have required floating-point arithmetic, is
  logged separately.
\end{itemize}

\textbf{5.5 Random-access metrics}

\begin{itemize}
\item
  Random-access wall-clock time $T(N, k)$ for
  $k \in \{1 \text{ KiB}, 4 \text{ KiB}, 16 \text{ KiB}, 64 \text{ KiB}\}$
  and $N$ sampled uniformly over the archive, with at least 1000
  samples per (archive size, $k$) pair. Report mean, median, 95th
  percentile.
\item
  Random-access throughput: number of independent random-access
  reads per second under serial query load.
\item
  Sync-point overhead: ratio cost as
  $(\text{archive size with sync points}) / (\text{archive size without sync points}) - 1$,
  reported at $K \in \{4, 16, 64, 256, 1024\}$ KiB.
\item
  Seek index size as a fraction of total archive size, with
  $W$-bounded versus non-$W$-bounded predictor variants separately
  reported.
\end{itemize}

\textbf{6. Statistical methodology}

\textbf{6.1 Sampling and runs}

For each (coder, corpus) pair, run the experiment 5 times with distinct
deterministic seeds, controlling for any non-determinism in baseline
coders (e.g., adaptive context mixing initialization). Report mean and
95\% confidence interval via bootstrap with 10000 resamples.

For randomized corpora (CommonCrawl subset, image samples), use
deterministic seeds for subset selection; verify subset hashes match
across runs.

\textbf{6.2 Effect-size reporting}

For comparisons between RNR variants and baselines, report:

\begin{itemize}
\item
  Mean rate difference in bits per byte (RNR minus baseline; negative is
  better).
\item
  Relative improvement as a percentage of the baseline rate.
\item
  95\% confidence interval on the difference, via paired bootstrap on
  per-block rates.
\item
  Cohen\textquotesingle s d effect size where applicable (block-level
  comparisons).
\end{itemize}

\textbf{6.3 Pre-registration}

To avoid p-hacking and result-cherry-picking, the experimental
hypothesis list (Section 2), the corpus list (Section 3), the metric
list (Section 5), and the analysis procedure (this section) are
committed in a pre-registration document with a publication timestamp,
prior to any experimental run. Pre-registration follows the protocol of
the OSF Registries or equivalent open-science platform.

\textbf{6.4 Multiple-comparisons correction}

With 9 hypotheses and multiple coder/corpus pairs each, the
per-comparison $\alpha$ of $0.05$ would inflate the family-wise
error rate substantially. Apply Holm-Bonferroni correction across all
primary-hypothesis tests. Secondary analyses (effect sizes,
exploratory measurements) report uncorrected confidence intervals
with a clear label distinguishing them from primary hypothesis tests.

\textbf{7. Reporting template}

Each experiment is reported in a uniform template for downstream
meta-analysis. The template is:
\begin{verbatim}
Experiment ID:        <slug, e.g., 'h2-pgdump-rnr-ii'>
Hypothesis:           <H1..H9 from Section 2>
Corpus:               <slug + SHA-256 hash>
Coder:                <RNR variant or baseline name + version>
Hardware:             <CPU, RAM, OS>
Software:             <library versions, RNR git commit hash>
Pre-registration ID:  <OSF URL or equivalent>
Runs:                 5

Result table:
  Rate (bpb):    mean +/- 95% CI
  Time (s/MiB):  mean +/- 95% CI
  Memory (MiB):  mean +/- 95% CI

Comparison to <baseline name>:
  Delta rate:                <value> bpb (<percentage>%)
  Bootstrap 95% CI on delta: [<low>, <high>]
  Cohen's d (block-level):   <value>

Hypothesis verdict:  <confirmed / falsified / inconclusive>
Notes:               <implementation details, anomalies, caveats>
\end{verbatim}

\textbf{8. Cross-platform reproducibility test (H8, H9)}

Theorems 10.1 and 10.2 reduce reproducibility to a conformance
specification. This section operationalizes the conformance test.

\textbf{8.1 Test corpus}

100 test blocks of varying size and content, drawn from the corpora in
Section 3. Each block is 1--16 KiB. The set is fixed by SHA-256 hashes
published with the test suite. The selection covers all RNR Types and at
least 10 blocks per Type.

\textbf{8.2 Test platforms}

\begin{itemize}
\item
  x86-64 Linux: Ubuntu 24.04 LTS, glibc 2.39, no AVX-512 (forced via
  taskset to a core without it).
\item
  x86-64 Linux with AVX-512 enabled (separate run, same OS).
\item
  ARM64 Linux: Apple Silicon under Asahi Linux, native code.
\item
  ARM64 macOS: Apple Silicon, native code, system clang/libc.
\item
  RISC-V via QEMU emulation (Ubuntu RISC-V image), no native hardware
  acceleration.
\item
  Windows x86-64 with MSVC runtime.
\end{itemize}

\textbf{8.3 Test procedure}

For each test block $X$ and each RNR variant, on each platform:
\begin{enumerate}
\item Run encoder on $X$ with deterministic seed; emit archive $A$.
\item Compute $\mathrm{SHA\text{-}256}(A)$.
\item Cross-tabulate SHA-256 values across platforms.
\item For all pairs $(\text{platform}_i, \text{platform}_j)$, record
$\text{byte-identity} = (\text{hash}_i = \text{hash}_j)$.
\item Run decoder on each $A$ on each platform; emit $\hat{X}$.
\item Verify $\hat{X} = X$ on every
$(\text{encoder platform}, \text{decoder platform})$ pair.
\end{enumerate}

Pass criterion: for every test block, all 6 platforms produce the same
archive hash, and every cross-decoded result equals the original. Fail
criterion: any byte difference. A single failure invalidates the
conformance claim for that RNR variant on that platform pair and
triggers diagnostic investigation (Section 8.4).

\textbf{8.4 Diagnostic procedure on failure}

If a platform pair produces divergent archives:

\begin{itemize}
\item
  Identify the earliest divergent byte in the archive.
\item
  Trace back to the encoder stage that produced it: predictor output,
  repair encoding, or metadata.
\item
  If predictor output diverges, check Theorem 10.1 conditions: integer
  bit-width, non-linearity table determinism, reduction order. Identify
  which condition is violated.
\item
  If repair encoding diverges, check entropy-coder state initialization
  and update order.
\item
  Publish the failure as an issue against the reference implementation,
  with the divergent input and the proposed fix.
\end{itemize}

\textbf{9. Intrinsic-dimension estimation procedure (H4)}

Remark 5.7 connects Theorem 5.6\textquotesingle s separation to
natural-data intrinsic dimension. This section operationalizes the
estimation.

\textbf{9.1 Embedding choice}

Intrinsic dimension is estimated in an embedding space, not raw
byte space, because raw bytes are unstructured. For each corpus,
compute byte $n$-gram embeddings at $n \in \{2, 4, 8, 16\}$. Each
block becomes a vector of counts over a vocabulary of distinct
$n$-grams observed in that corpus. Vocabulary is capped at 65536 by
frequency.

Alternative embedding: use the byte-level transformer\textquotesingle s
penultimate-layer hidden state at each position, averaged over the
block. This requires running the RNR predictor in inference mode, which
is acceptable since the predictor is already part of the framework.

\textbf{9.2 Estimator}

Apply the Levina-Bickel maximum-likelihood estimator (Levina-Bickel
2004) with $k$ nearest neighbors. For each sample point $x_i$ and
each of its $k$ nearest neighbors $x_{i,1}, \ldots, x_{i,k}$ (by
Euclidean distance in the embedding space):
\[
\begin{aligned}
T_k(x_i)         &= \text{distance to $k$-th nearest neighbor,} \\
\hat{d}(x_i, k)  &= \frac{k - 1}{\sum_{j=1}^{k-1} \log(T_k(x_i) / T_j(x_i))}, \\
\hat{d}(k)       &= \text{mean over all sample points of } \hat{d}(x_i, k).
\end{aligned}
\]
Report $\hat{d}$ for $k \in \{5, 10, 20, 50\}$ and stable across
these $k$ values. If $\hat{d}$ varies substantially across $k$, the
corpus does not exhibit a clean low-dimensional manifold structure.

\textbf{9.3 Validation against alternative estimators}

Cross-check with at least one alternative:

\begin{itemize}
\item
  Correlation dimension via the Grassberger-Procaccia algorithm.
\item
  Persistent homology dimension via the persistent-homology approach
  (Bubenik 2015 or similar).
\item
  Two-NN dimension estimator (Facco et al. 2017).
\end{itemize}

Report agreement or disagreement among estimators. Strong
disagreement (factor of 2 or more) indicates the corpus may not
admit a manifold description; in that case Remark 5.7's
quantification of $\mathrm{TC}$ does not transfer cleanly to that
corpus.

\textbf{10. Suggested execution order}

The experimental program has dependencies: reproducibility tests must
pass before rate comparisons are meaningful, intrinsic-dimension
estimates must precede the H3/H4 analysis, etc. The suggested order is:

\emph{\textbf{Phase 1: implementation and reproducibility (weeks 1--3)}}

\begin{itemize}
\item
  Build reference implementations of all RNR variants meeting Theorem
  10.1 conditions.
\item
  Run H8 conformance test. Iterate until all platforms produce
  byte-identical archives.
\item
  Build adaptive-mode variants with compatible optimizers and run H9.
\end{itemize}

\emph{\textbf{Phase 2: information-theoretic measurements (weeks 4--6)}}

\begin{itemize}
\item
  Estimate $\mathrm{TC}(D_N)$ and intrinsic dimension $d$ for each
  corpus (H3, H4).
\item
  Publish the estimates as a standalone artifact (pre-registration
  permits early publication of corpus characterization).
\item
  Estimate $I(H ; L \mid C)$ for H6.
\end{itemize}

\emph{\textbf{Phase 3: rate measurements (weeks 7--12)}}

\begin{itemize}
\item
  Run all coder/corpus pairs at fixed predictor capacity. Apply the
  reporting template from Section 7.
\item
  H1 on text and code corpora.
\item
  H2 on structured binary corpora.
\item
  H5 on heterogeneous corpora (after composition implementation is
  complete).
\item
  H6 cross-checks (rate gain in Type-III-B against measured
  $I(H ; L \mid C)$).
\item
  H7 on image corpora (bits-back over VAE).
\item
  H10 and H11 random-access scaling tests, run after H1 through H7
  confirm baseline encoder correctness. Cover at least three archive
  sizes (100 MiB, 10 GiB, 100 GiB) with 1000 random-access samples each.
\end{itemize}

\emph{\textbf{Phase 4: analysis and reporting (weeks 13--15)}}

\begin{itemize}
\item
  Apply Holm-Bonferroni correction across all primary hypothesis tests.
\item
  Compile per-hypothesis verdicts and aggregate effect sizes.
\item
  Draft companion experimental paper for publication, with all data and
  code released as open artifacts.
\end{itemize}

Total estimated duration: 15 weeks for one researcher, less with
parallelism on implementation. The reproducibility tests in Phase 1 are
the riskiest stage; they may take longer if conformance issues are
subtle. Phases 2 and 3 run mostly automated once implementations are
stable.

\textbf{11. Risks and contingencies}

Anticipated risks and contingency plans:

\begin{itemize}
\item
  Reproducibility may fail on RISC-V or Windows due to subtle
  integer-arithmetic differences (e.g., division rounding modes, shift
  behavior). Contingency: replace the failing operations with portable
  alternatives explicitly listed in the model description; if not
  feasible, restrict the conformance claim to a subset of platforms and
  document this in the failure analysis.
\item
  Predictor training may not converge under integer-arithmetic
  constraints. Contingency: train in floating point, then quantize to
  integer; verify post-quantization that the model still satisfies
  Theorem 10.1 conditions. Accept some predictor quality loss in
  exchange for reproducibility.
\item
  Manifold-hypothesis estimators may give inconsistent intrinsic
  dimensions on certain corpora. Contingency: report the disagreement
  honestly and label the corresponding corpus as outside Remark
  5.7\textquotesingle s predicted regime. The
  framework\textquotesingle s conditional theorems are unaffected; only
  the empirical applicability for that corpus is in question.
\item
  Bits-back implementation may have variational-training instabilities
  (typical for VAE-style training). Contingency: use a published
  Bit-Swap or HiLLoC reference implementation as the RNR-III-C-V
  baseline rather than training from scratch; this measures the
  framework integration, not the variational fit.
\item
  Predictor cost L(M) may dominate rate at small block sizes, masking
  the structural advantages of each Type. Contingency: report rates
  separately at multiple total-archive sizes (10 MB, 100 MB, 1 GB) so
  amortization is visible.
\end{itemize}

\textbf{12. Open artifacts}

To support independent verification, the experimental program commits to
publishing:

\begin{itemize}
\item
  All RNR reference implementations under a permissive license.
\item
  Pre-registration document with hypothesis list and analysis plan,
  timestamped before any run.
\item
  Corpus manifests with SHA-256 hashes and provenance URLs.
\item
  All raw measurement outputs (per-block rates, timings, memory
  profiles) in a machine-readable format.
\item
  Statistical analysis scripts used to produce the reported confidence
  intervals.
\item
  Final paper drafts, including any negative or inconclusive results,
  even those that falsify framework hypotheses.
\end{itemize}

Failed-hypothesis publication is a hard commitment: a result that
falsifies H1, H2, H5, H6, or H7 is published with the same rigor as a
confirming result. The framework\textquotesingle s theorems are
unaffected by such falsification; what they show is which conjectures
about practical realization do not hold. This information is more
valuable than another confirming paper because it identifies where the
framework\textquotesingle s predictions diverge from natural-data
reality.

\textbf{13. Closing}

This document is an executable scaffolding, not a research result.
Following it produces measurements that resolve every empirical
conjecture made by the main RNR Coding paper. The combination of
pre-registration, open artifacts, and explicit falsification criteria
makes the resulting experimental paper falsifiable by independent
replication. Whatever the outcome, the framework\textquotesingle s
theoretical content (the 17 numbered results in v0.5) is unchanged; the
empirical paper says only how well the framework\textquotesingle s
conditional predictions match natural-data reality.

Two ways to read this document. As a research plan: it tells you what to
build, what to measure, and what to publish, in what order. As a
defense: it shows a reviewer that the main paper\textquotesingle s
empirical conjectures are accompanied by a concrete falsification
procedure, which is the only honest way to handle theory-only papers
that make quantitative predictions about real data.

\end{document}
